\documentclass{beamer}
\usepackage[en]{zafu}

% PDF metadata and safe bookmark settings (no visual changes)
\hypersetup{
  pdftitle={An Informer-Based Calibration-Free LIBS Spectral Identification},
  pdfauthor={Birrul Walidain; Nasrullah Idris; Khairun Saddami; Natasya Yuza},
  pdfsubject={LIBS; Informer; Calibration-free; Spectral Analysis},
  pdfkeywords={LIBS, Informer, Spectral Identification}
}
% Avoid warnings from superscripts and line breaks in PDF strings
\pdfstringdefDisableCommands{%%
  \def\\{}%% ignore manual line breaks in bookmarks
  \def\textsuperscript#1{#1}%% drop superscript styling in bookmarks
}

% Metadata (7-minute presentation)
\title[Calibration-Free LIBS w/ Informer]{An Informer-Based Calibration-Free LIBS Spectral Identification}
\subtitle{Case Study: Aceh Traditional Women's Medicine}
\author[Walidain et al.]{
  Birrul Walidain\textsuperscript{1}, Nasrullah Idris\textsuperscript{1*}, \\
  Khairun Saddami\textsuperscript{2}, Natasya Yuza\textsuperscript{1}
}
\institute{
  \small
  \textsuperscript{1}Department of Physics, Faculty of Mathematics and Natural Sciences\\
  \textsuperscript{2}Department of Electrical and Computer Engineering\\
  \textbf{Universitas Syiah Kuala}, Banda Aceh, Indonesia, 23111\\
  \vspace{0.5em}
  \textsuperscript{*}Corresponding author: \texttt{birrul@mhs.usk.ac.id}
}
\date{September 5, 2025}

% Optional: small logo on title (uses your USK logo)
\titlegraphic{\includegraphics[width=0.16\textwidth]{src/logo-usk.png}}

\begin{document}

% Slide 1: Title Slide (0:00–0:30)
\maketitle

\section{Introduction}

% Slide 2: Research Objectives (0:30–1:15)
\begin{frame}{Research Objectives}
  \begin{block}{Main Goal}
    Develop calibration-free LIBS identification using Informer deep learning with synthetic physics-based training
  \end{block}
  
  \vspace{1em}
  
  \begin{itemize}
    \item \textbf{Innovation}: First use of Informer architecture for LIBS spectral analysis
    \item \textbf{Physics approach}: Saha-Boltzmann equations + NIST database simulation
    \item \textbf{Efficiency gain}: ProbSparse attention reduces complexity to $\mathcal{O}(L \log L)$
    \item \textbf{Validation}: Synthetic-trained model tested on real experimental data
  \end{itemize}
  
  \vspace{0.5em}
  \begin{alertblock}{Research Question}
    Can synthetic LIBS spectra based on fundamental plasma physics replace expensive calibration procedures?
  \end{alertblock}
\end{frame}

% Slide 3: Problem Statement (1:15–1:45)
\begin{frame}{Why LIBS Analysis is Challenging}
  \begin{columns}[T]
    \begin{column}{0.48\textwidth}
      \textbf{Spectral Complexity:}
      \begin{itemize}
        \item Al/Ca interference at 393-396 nm
        \item Ion-neutral equilibrium shifts
        \item Temperature-dependent dominance
        \item Multi-element overlapping
      \end{itemize}
    \end{column}
    \begin{column}{0.48\textwidth}
      \textbf{Analysis Challenges:}
      \begin{itemize}
        \item Ca II (393.37 nm) vs Al I overlap
        \item Saha equilibrium complexity
        \item Matrix effect variations
        \item Manual peak deconvolution
      \end{itemize}
    \end{column}
  \end{columns}
  
  \vspace{0.5em}
  \begin{center}
    \includegraphics[width=0.8\textwidth]{images/Al-Ca-Fe-Si_T8000K_ne1e17_wl390-400nm.png}
  \end{center}
  
  \vspace{1em}
  \begin{alertblock}{Real Example from Research}
    At T=8000K: Ca II dominates (100\% intensity), Al I weak (14.6\% Saha fraction)\\
    Traditional methods struggle with such dynamic equilibria
  \end{alertblock}
\end{frame}

\section{Methodology}

% Slide 4: Informer Architecture (1:45–2:30)
\begin{frame}{Informer-Based Solution}
  \begin{block}{Key Innovation: ProbSparse Self-Attention}
    Reduces complexity from $\mathcal{O}(L^2)$ to $\mathcal{O}(L \log L)$ for efficient spectral processing
  \end{block}
  
  \vspace{0.5em}
  
  \begin{columns}[T]
    \begin{column}{0.48\textwidth}
      \textbf{Technical Implementation:}
      \begin{itemize}
        \item 2-layer encoder (optimal depth)
        \item MultiLabelFocalLoss for class imbalance
        \item Automatic Mixed Precision (AMP)
        \item 18-element multi-label classification
      \end{itemize}
      
      \vspace{0.5em}
      \begin{center}
        \includegraphics[width=0.9\textwidth]{images/3-Arsitektur-Informer.drawio.pdf}
      \end{center}
    \end{column}
    \begin{column}{0.48\textwidth}
      \textbf{Training Details:}
      \begin{itemize}
        \item 70 epochs, stable convergence
        \item Focal Loss $\alpha$=0.25, $\gamma$=2.0
        \item Dropout regularization
        \item Physics-guided attention
      \end{itemize}
      
      \vspace{0.5em}
      \begin{center}
        \includegraphics[width=0.9\textwidth]{images/training_history_plot.png}
      \end{center}
    \end{column}
  \end{columns}
  
  \vspace{0.5em}
  \begin{exampleblock}{Key Advantage}
    ProbSparse Self-Attention efficiently handles high-resolution LIBS spectral sequences
  \end{exampleblock}
\end{frame}

% Slide 5: Dataset and Training (2:30–3:15)
\begin{frame}{Physics-Based Synthetic Training}
  \begin{block}{Synthetic Spectra Generation using NIST Data}
    Saha-Boltzmann equations + NIST Atomic Spectra Database (ASD) transitions
  \end{block}
  
  \begin{columns}[T]
    \begin{column}{0.48\textwidth}
      \textbf{Physics Implementation:}
      \begin{itemize}
        \item NIST ASD atomic transition data
        \item Saha equilibrium calculations
        \item Boltzmann population statistics
        \item Gaussian line profile broadening
      \end{itemize}
      
      \vspace{0.5em}
      \begin{center}
        \includegraphics[width=0.9\textwidth]{images/Al-Ca-Fe-Si_T8000K_ne1e17-NIST.png}
      \end{center}
    \end{column}
    \begin{column}{0.48\textwidth}
      \textbf{Simulation Features:}
      \begin{itemize}
        \item Local Thermal Equilibrium (LTE)
        \item Controlled plasma conditions
        \item Multi-element interference modeling
        \item Realistic spectral complexity
      \end{itemize}
      
      \vspace{0.5em}
      \begin{center}
        \includegraphics[width=0.9\textwidth]{images/Simulasi.drawio.pdf}
      \end{center}
    \end{column}
  \end{columns}
\end{frame}

\section{Results}

% Slide 6: Performance Results (3:15–4:15)
\begin{frame}{Experimental Results}
  \begin{columns}[T]
    \begin{column}{0.48\textwidth}
      \textbf{Model Performance:}
      \begin{itemize}
        \item \textcolor{green}{\textbf{F1-Score: 0.6445}} (2-layer optimal)
        \item \textbf{Precision:} 0.5938 (macro avg)
        \item \textbf{Recall:} 0.7733 (high sensitivity)
        \item 3-layer: 0.5199 (overfitting)
      \end{itemize}
    \end{column}
    \begin{column}{0.48\textwidth}
      \textbf{Element-Specific Results:}
      \begin{itemize}
        \item Ca: F1=0.8065 (excellent)
        \item Si: F1=0.8397 (excellent)
        \item Al: F1=0.6665 (good)
        \item Fe: F1=0.1761 (challenging)
      \end{itemize}
    \end{column}
  \end{columns}
  
  \vspace{1em}
  
  \begin{block}{Key Findings}
    High recall (0.77) shows excellent sensitivity; Lower precision indicates interference challenges - consistent with physics expectations
  \end{block}
\end{frame}

% Slide 7: Case Study Application (4:15–5:15)
\begin{frame}{Traditional Medicine Analysis}
  \begin{block}{Real-World Application: Aceh Traditional Women's Medicine}
    Practical demonstration of calibration-free spectral identification on complex herbal matrices
  \end{block}
  
  \begin{columns}[T]
    \begin{column}{0.48\textwidth}
      \textbf{Elements Identified:}
      \begin{itemize}
        \item \textbf{Major organic:} C, N, O
        \item \textbf{Nutritional minerals:} Ca, Mg
        \item \textbf{Trace elements:} Al, Fe, Mn, Na
        \item Complex matrix deconvolution
      \end{itemize}
    \end{column}
    \begin{column}{0.48\textwidth}
      \textbf{Key Achievements:}
      \begin{itemize}
        \item Detailed spectral deconvolution
        \item No experimental calibration needed
        \item Multi-element interference handled
        \item Practical screening capability
      \end{itemize}
    \end{column}
  \end{columns}
  
  \vspace{1em}
  \begin{exampleblock}{Clinical Impact}
    Rapid screening tool for traditional medicine quality control and safety assessment
  \end{exampleblock}
\end{frame}

\section{Conclusion}

% Slide 8: Conclusions and Future Work (5:15–6:15)
\begin{frame}{Conclusions \& Future Directions}
  \begin{block}{Key Achievements}
    \begin{itemize}
      \item \textbf{Synthetic-only training}: Calibration-free approach using NIST-based synthetic spectra
      \item \textbf{Practical performance}: F1=0.6445 with optimized 2-layer Informer architecture
      \item \textbf{Real-world application}: Traditional medicine elemental analysis demonstrated
    \end{itemize}
  \end{block}
  
  \vspace{1em}
  
  \begin{columns}[T]
    \begin{column}{0.48\textwidth}
      \textbf{Research Contributions:}
      \begin{itemize}
        \item ProbSparse Self-Attention for LIBS
        \item Physics-based synthetic training
        \item Multi-element deconvolution capability
      \end{itemize}
    \end{column}
    \begin{column}{0.48\textwidth}
      \textbf{Future Research:}
      \begin{itemize}
        \item Extended element database
        \item Quantitative concentration analysis  
        \item Non-LTE plasma conditions
      \end{itemize}
    \end{column}
  \end{columns}
  
  \vspace{1em}
  \begin{alertblock}{Impact}
    Bridges advanced computational models with public health needs for rapid screening
  \end{alertblock}
\end{frame}

% Slide 9: Q&A and Acknowledgments (6:15–7:00 + 3 min Q&A)
\begin{frame}{}
  \vfill
  \begin{center}
    {\Huge \textcolor{uskGold}{\textbf{Thank You!}}}
    
    \vspace{1em}
    
    {\Large Questions \& Discussion}
    
    \vspace{1.5em}
    
    {\large \textcolor{uskGold}{\textbf{Contacts:}}}
    
    \vspace{0.5em}
    
    {\normalsize 
    \textbf{Birrul Walidain:} \texttt{birrul@mhs.usk.ac.id}\\
    \textbf{Nasrullah Idris:} \texttt{nasrullah.idris@usk.ac.id}
    }
  \end{center}
  \vfill
  
  \begin{center}
    {\footnotesize \textcolor{gray}{\textbf{Acknowledgments}}}
    
    \vspace{0.3em}
    
    {\scriptsize \textcolor{gray}{
    This research was supported by BRIN High Performance Computing\\
    (Proposals 224246 \& 189541) and LPPM USK (Patent Application S00202506246).\\
    Thanks to AIC-SE USK organizing committee.
    }}
  \end{center}
\end{frame}

% Slide 10: Template Credits
\begin{frame}{}
  \vfill
  \begin{center}
    {\Large \textcolor{uskGold}{\textbf{Template Credits}}}
    
    \vspace{1.5em}
    
    {\large Presentation template based on:}
    
    \vspace{1em}
    
    {\LARGE \textbf{ZAFU Beamer Theme}}\\
    \vspace{0.5em}
    {\normalsize Original by: \textcolor{gray}{github.com/ke1ewang/ZAFU-Beamer-Theme}}
    
    \vspace{1.5em}
    
    {\large \textcolor{uskGold}{\textbf{USK Adaptation}}}\\
    \vspace{0.5em}
    {\normalsize Adapted \& Modified by: \textbf{Birrul Walidain}}\\
    {\small \textcolor{gray}{Department of Physics, USK}}
    
    \vspace{1em}
    
    {\small \textcolor{gray}{September 2025}}
  \end{center}
  \vfill
\end{frame}

\end{document}
